\section{Garantia de Minimalidade}
\label{sec:minimizacao}

Um autômato é \emph{mínimo} quando nenhum de seus estados pode ser
removido ou fundido com outro sem alterar o conjunto de placas
aceitas. Mostrar que o modelo é mínimo garante que ele não tem
etapas redundantes: cada estado guarda uma informação que nenhum
outro guarda.

\subsection{Método}

Dois estados $p$ e $q$ são \emph{distinguíveis} quando existe uma
cadeia $w$ tal que exatamente um entre $\hat{\delta}(p,w)$ e
$\hat{\delta}(q,w)$ é final. Estados que não podem ser
distinguidos são \emph{equivalentes} e podem ser fundidos
\cite{hopcroft2002}.

O algoritmo de refinamento de partições (Moore) funciona assim:
\begin{enumerate}
  \item $P_0 = \{\,F,\; Q \setminus F\,\}$.
  \item Em cada rodada, dois estados continuam no mesmo bloco
        apenas se, para cada classe de símbolos, suas transições
        levam a estados que estavam em um mesmo bloco na rodada
        anterior.
  \item Repete-se até que a partição não mude. Os blocos finais
        são os estados do autômato mínimo.
\end{enumerate}

Como os símbolos de uma mesma classe se comportam do mesmo modo,
basta comparar as transições por classe ($\Lt$, $\Nm$ e $\Hf$). O
método exige que todos os estados sejam alcançáveis a partir de
$q_0$, o que vale para todos os autômatos deste trabalho.


\subsection{AFD 01}

A partição inicial separa o único estado final dos demais:
\[
  P_0 = \bigl\{\, \{q_8\},\; \{q_0,q_1,q_2,q_3,q_4,q_5,q_6,q_7,q_R\} \,\bigr\}.
\]

\begin{table}[ht]
\centering
\small
\renewcommand{\arraystretch}{1.3}
\begin{tabular}{c p{3.4cm} p{7.6cm}}
\toprule
Rodada & Blocos isolados & Motivo \\
\midrule
$P_1$ & $\{q_7\}$
  & $q_7 \xrightarrow{\Nm} q_8$, que é final; nenhum outro estado do bloco chega a um final em um passo. \\
$P_2$ & $\{q_6\}$
  & $q_6 \xrightarrow{\Nm} q_7$, já isolado. \\
$P_3$ & $\{q_5\}$
  & $q_5 \xrightarrow{\Nm} q_6$, já isolado. \\
$P_4$ & $\{q_3,q_4\}$ separado de $\{q_0,q_1,q_2,q_R\}$
  & $q_3 \xrightarrow{\Nm} q_5$ e $q_4 \xrightarrow{\Nm} q_5$, bloco já isolado; os demais não. \\
$P_5$ & $\{q_3\}$, $\{q_4\}$, $\{q_2\}$
  & $q_3 \xrightarrow{\Hf} q_4$, mas $q_4 \xrightarrow{\Hf} q_R$ (outro bloco); $q_2 \xrightarrow{\Lt} q_3$. \\
$P_6$ & $\{q_1\}$
  & $q_1 \xrightarrow{\Lt} q_2$, já isolado. \\
$P_7$ & $\{q_0\}$, $\{q_R\}$
  & $q_0 \xrightarrow{\Lt} q_1$, enquanto $q_R \xrightarrow{\Lt} q_R$. \\
\bottomrule
\end{tabular}
\caption{Refinamento de partições do AFD 01.}
\label{tab:min1}
\end{table}

Na rodada $P_7$ todos os blocos têm um único estado, e nenhuma nova
separação é possível. Portanto os 10 estados do AFD 01 são
distinguíveis dois a dois, e o autômato é mínimo. O par que exigiu
mais cuidado foi $q_3$ e $q_4$: ambos aceitam \texttt{1234} como
continuação, mas apenas $q_3$ aceita \texttt{-1234}. Outro exemplo:
a partir de $q_0$ a cadeia \texttt{ABC1234} é aceita, e a partir de
$q_R$ nenhuma cadeia é aceita.


\subsection{AFD 02 e AFD 03}

No AFD 02, a partição inicial é
$P_0 = \bigl\{ \{q_7\},\; \{q_0,\dots,q_6,q_R\} \bigr\}$. Como o
autômato é uma cadeia de etapas, cada rodada isola o estado
anterior da cadeia.

\begin{table}[ht]
\centering
\small
\renewcommand{\arraystretch}{1.3}
\begin{tabular}{c p{3.4cm} p{7.6cm}}
\toprule
Rodada & Bloco isolado & Motivo \\
\midrule
$P_1$ & $\{q_6\}$ & $q_6 \xrightarrow{\Nm} q_7$, que é final. \\
$P_2$ & $\{q_5\}$ & $q_5 \xrightarrow{\Nm} q_6$, já isolado. \\
$P_3$ & $\{q_4\}$ & $q_4 \xrightarrow{\Lt} q_5$, já isolado. \\
$P_4$ & $\{q_3\}$ & $q_3 \xrightarrow{\Nm} q_4$, já isolado. \\
$P_5$ & $\{q_2\}$ & $q_2 \xrightarrow{\Lt} q_3$, já isolado. \\
$P_6$ & $\{q_1\}$ & $q_1 \xrightarrow{\Lt} q_2$, já isolado. \\
$P_7$ & $\{q_0\}$, $\{q_R\}$ & $q_0 \xrightarrow{\Lt} q_1$, enquanto $q_R \xrightarrow{\Lt} q_R$. \\
\bottomrule
\end{tabular}
\caption{Refinamento de partições do AFD 02.}
\label{tab:min2}
\end{table}

Todos os blocos de $P_7$ têm um único estado, logo os 9 estados do
AFD 02 são distinguíveis e o autômato é mínimo.

O AFD 03 tem a mesma estrutura de cadeia. A única diferença é que
$q_4$ separa as letras em $\LAJ$ e $\Lt \setminus \LAJ$, e isso
não altera a sequência de rodadas. Assim, o AFD 03 também é
mínimo, com 9 estados.

\subsection{Validador combinado em um único autômato}

Para ter um único autômato que decide a validade, constrói-se o
\emph{produto} dos AFDs 01 e 02. Seus estados são pares
(estado do AFD 01, estado do AFD 02), e um par é final se algum dos
componentes é final. Apenas 14 pares são alcançáveis.

\begin{table}[ht]
\centering
\small
\renewcommand{\arraystretch}{1.15}
\begin{tabular}{cc@{\hspace{2em}}cc}
\toprule
Estado & Par & Estado & Par \\
\midrule
$u_0$ & $(q_0,q_0)$ & $u_7$ & $(q_R,q_5)$ \\
$u_1$ & $(q_1,q_1)$ & $u_8$ & $(q_6,q_R)$ \\
$u_2$ & $(q_2,q_2)$ & $u_9$ & $(q_R,q_6)$ \\
$u_3$ & $(q_3,q_3)$ & $u_{10}$ & $(q_7,q_R)$ \\
$u_4$ & $(q_4,q_R)$ & $u_{11}$ (final) & $(q_R,q_7)$ \\
$u_5$ & $(q_5,q_4)$ & $u_{12}$ (final) & $(q_8,q_R)$ \\
$u_6$ & $(q_5,q_R)$ & $u_{13}$ & $(q_R,q_R)$ \\
\bottomrule
\end{tabular}
\caption{Estados alcançáveis do produto.}
\label{tab:produto}
\end{table}

\begin{table}[ht]
\centering
\small
\renewcommand{\arraystretch}{1.3}
\begin{tabular}{c p{3.4cm} p{7.6cm}}
\toprule
Rodada & Blocos isolados & Motivo \\
\midrule
$P_0$ & $\{u_{11},u_{12}\}$ & estados finais. \\
$P_1$ & $\{u_9,u_{10}\}$ & $u_9 \xrightarrow{\Nm} u_{11}$ e $u_{10} \xrightarrow{\Nm} u_{12}$, ambos finais. \\
$P_2$ & $\{u_7,u_8\}$ & $u_7 \xrightarrow{\Nm} u_9$ e $u_8 \xrightarrow{\Nm} u_{10}$. \\
$P_3$ & $\{u_5\}$, $\{u_6\}$ & $u_5 \xrightarrow{\Lt} u_7$, mas $u_6 \xrightarrow{\Lt} u_{13}$; ambos chegam ao bloco $\{u_7,u_8\}$ com $\Nm$. \\
$P_4$ & $\{u_3\}$, $\{u_4\}$ & $u_3 \xrightarrow{\Nm} u_5$ e $u_4 \xrightarrow{\Nm} u_6$. \\
$P_5$ & $\{u_2\}$ & $u_2 \xrightarrow{\Lt} u_3$. \\
$P_6$ & $\{u_1\}$ & $u_1 \xrightarrow{\Lt} u_2$. \\
$P_7$ & $\{u_0\}$, $\{u_{13}\}$ & $u_0 \xrightarrow{\Lt} u_1$, enquanto $u_{13} \xrightarrow{\Lt} u_{13}$. \\
\bottomrule
\end{tabular}
\caption{Refinamento de partições do produto.}
\label{tab:min3}
\end{table}

A partição final tem 11 blocos. Três pares de estados são
equivalentes: $\{u_7,u_8\}$, $\{u_9,u_{10}\}$ e $\{u_{11},u_{12}\}$.
Em palavras: $u_7$ (Mercosul) e $u_8$ (padrão anterior) esperam
exatamente os mesmos dois algarismos finais; $u_9$ e $u_{10}$
esperam um; $u_{11}$ e $u_{12}$ já terminaram. O que vem depois é
igual, então o autômato mínimo os trata como um só estado.

Há um custo nessa fusão: o estado final único deixa de informar
\emph{qual} padrão foi reconhecido. Se o sistema precisar
classificar a placa, os dois finais devem ser mantidos separados na
partição inicial, e nesse caso nenhuma fusão ocorre e o autômato
fica com 14 estados.


\subsection{Resumo}

\begin{table}[ht]
\centering
\renewcommand{\arraystretch}{1.2}
\begin{tabular}{lccl}
\toprule
Autômato & Estados antes & Estados depois & Conclusão \\
\midrule
AFD 01 (anterior)     & 10 & 10 & já é mínimo \\
AFD 02 (Mercosul)     & 9  & 9  & já é mínimo \\
AFD 03 (convertido)   & 9  & 9  & já é mínimo \\
Produto (validador)   & 14 & 11 & reduzido por 3 fusões \\
\bottomrule
\end{tabular}
\caption{Resultado da minimização.}
\label{tab:resumo-min}
\end{table}